% Discussion for item 3.
% scripts/make_latex.py will NEVER overwrite this file -- edit it freely.

Other programs running on the system really do reduce how many times we can
fork. On its own the program managed \textbf{4295} successful \texttt{fork()}
calls; with 200 extra \texttt{sleep} processes running it managed \textbf{4226};
after those were killed it was back to \textbf{4296}. Measurements 1 and 3 differ
by only 1, so the drop of 69 is well outside the run-to-run variation.

The reason is that whatever runs out is not owned by our program. It is shared
across everything the user is running, so processes that have nothing to do with
us --- shells, editors, those 200 \texttt{sleep} processes --- consume part of the
same budget and leave less for us. Killing them gives it straight back, which is
why the third measurement returns to the first.

How much less depends on which resource is the scarce one. Here the drop was 69
rather than 200, because the limit being hit was memory and not a limit on the
\emph{number} of processes: the 200 \texttt{sleep} processes consumed about
256~MB between them (free memory before each run: 8269~MB, then 8013~MB), and at
this depth each of our own processes costs several megabytes, so 256~MB is worth
roughly 69 of them. Had the per-user process quota been the binding limit
instead, 200 processes would have cost exactly 200 forks, one seat each.

It is worth being clear about what stopped the run, because it was not the
process limit --- this machine allows 35\,595 processes per user and
\texttt{fork()} never returned an error --- nor a shortage of PIDs, of which there
are 4\,194\,304. Memory ran out first, since each process costs the kernel a
process table entry, a kernel stack and page tables. The counts above are
therefore a floor rather than a true ceiling: each run was stopped once free
memory fell below 800~MB, so that the OOM killer was never left to pick a victim
of its own.
